
% ==========================================================================
% cc-latex.tex — Shared review and collaboration macros for Overleaf
% ==========================================================================
%
% SETUP:
%   1. Upload this file to your Overleaf project
%   2. Add 
% ==========================================================================
% cc-latex.tex — Shared review and collaboration macros for Overleaf
% ==========================================================================
%
% SETUP:
%   1. Upload this file to your Overleaf project
%   2. Add 
% ==========================================================================
% cc-latex.tex — Shared review and collaboration macros for Overleaf
% ==========================================================================
%
% SETUP:
%   1. Upload this file to your Overleaf project
%   2. Add 
% ==========================================================================
% cc-latex.tex — Shared review and collaboration macros for Overleaf
% ==========================================================================
%
% SETUP:
%   1. Upload this file to your Overleaf project
%   2. Add \input{cc-latex} in your preamble (after \usepackage{xcolor})
%   3. Make sure xcolor and changes packages are available
%
% TOGGLES (set to true/false in this file to control visibility):
%   \COMMENTStrue / \COMMENTSfalse    — author comments (\simo, \bob, etc.)
%   \coloringtrue / \coloringfalse    — revision coloring (\rone, \rtwo)
%   \TODOStrue / \TODOSfalse          — todo markers (\todo)
%
% QUICK REFERENCE:
%   \red{text}                  — always-on red text (not toggleable)
%   \todo{remember to fix this} — red TODO text (toggle: todos)
%   \simo{text}{comment}        — inline author comment (toggle: COMMENTS)
%   \rone{revised text}          — blue revision marker (toggle: coloring)
%   \rtwo{revised text}         — violet revision marker (toggle: coloring)
%   \added{text}  \deleted{text}  \replaced{new}{old}  — changes package
%
% FOR FINAL/CAMERA-READY: set all three toggles to false, and switch
%   \usepackage[draft]{changes} to \usepackage[final]{changes}
% ==========================================================================


% --------------------------------------------------------------------------
% UTILITY COMMANDS
% --------------------------------------------------------------------------

% Paint text red — handy for quick emphasis or flagging
% Usage: \red{some important text}
\newcommand{\red}[1]{\textcolor{red}{#1}}


% --------------------------------------------------------------------------
% TODO MARKERS
% --------------------------------------------------------------------------
% Toggle: \TODOStrue / \TODOSfalse

\newif\ifTODOS
\TODOStrue

\newcommand{\todo}[1]{%
  \ifTODOS\textcolor{red}{[TODO: #1]}\fi%
}


% --------------------------------------------------------------------------
% AUTHOR COMMENT COMMANDS
% --------------------------------------------------------------------------
% Toggle: set \COMMENTStrue or \COMMENTSfalse

\newif\ifCOMMENTS
\COMMENTStrue

\ifCOMMENTS
\newcommand{\simo}[2]{\textcolor{red}{#1\bf{ (Simo: #2)}}}
\newcommand{\bob}[2]{\textcolor{blue}{#1\bf{ (Bob: #2)}}}
\newcommand{\alice}[2]{\textcolor{teal}{#1\bf{ (Alice: #2)}}}
\newcommand{\negatron}[2]{\textcolor{purple}{#1\bf{ (Negatron: #2)}}}
\else
\newcommand{\simo}[2]{#1}
\newcommand{\bob}[2]{#1}
\newcommand{\alice}[2]{#1}
\newcommand{\negatron}[2]{#1}
\fi

% Usage: \simo{test sentence}{maybe rephrase this part}
%        \bob{some text}{I disagree with this}
%        \alice{a claim}{needs citation}
%        \negatron{paragraph}{this section is weak}
%
% TO ADD A NEW AUTHOR:
%   1. Pick a color (e.g. orange, brown, cyan — see xcolor docs)
%   2. Add inside BOTH the \ifCOMMENTS and \else branches:
%      \ifCOMMENTS block:
%        \newcommand{\yourname}[2]{\textcolor{yourcolor}{#1\bf{ (Yourname: #2)}}}
%      \else block:
%        \newcommand{\yourname}[2]{#1}
%   3. Optionally add a \definechangesauthor line below for the changes package


% --------------------------------------------------------------------------
% CHANGE TRACKING (changes package)
% --------------------------------------------------------------------------
% https://mirrors.ibiblio.org/CTAN/macros/latex/contrib/changes/changes.english.pdf
\usepackage[draft]{changes}  % shows changes
%\usepackage[final]{changes}  % hides changes

% Define authors for the changes package
\definechangesauthor[name={Simo}, color=red]{SH}
\definechangesauthor[name={Bob}, color=blue]{BO}
\definechangesauthor[name={Alice}, color=teal]{AL}
\definechangesauthor[name={Negatron}, color=purple]{NE}

% Basic usage:
%   \added{This is added text.}
%   \deleted{This text is deleted.}
%   \replaced{new text}{old text}
%
% With author tag:
%   \added[SH]{Author-tagged addition.}
%   \deleted[BO]{Author-tagged deletion.}
%   \replaced[AL]{newer text}{older text}
% --------------------------------------------------------------------------


% --------------------------------------------------------------------------
% REVISION COLORING
% --------------------------------------------------------------------------
% Toggle: \coloringtrue / \coloringfalse
% Use \rone{} for round 1 revisions, \rtwo{} for round 2

\newif\ifcoloring
\coloringtrue

\newcommand{\rone}[1]{%
  \ifcoloring{\color{blue}#1}\else#1\fi%
}

\newcommand{\rtwo}[1]{%
  \ifcoloring{\color{violet}#1}\else#1\fi%
}
 in your preamble (after \usepackage{xcolor})
%   3. Make sure xcolor and changes packages are available
%
% TOGGLES (set to true/false in this file to control visibility):
%   \COMMENTStrue / \COMMENTSfalse    — author comments (\simo, \bob, etc.)
%   \coloringtrue / \coloringfalse    — revision coloring (\rone, \rtwo)
%   \TODOStrue / \TODOSfalse          — todo markers (\todo)
%
% QUICK REFERENCE:
%   \red{text}                  — always-on red text (not toggleable)
%   \todo{remember to fix this} — red TODO text (toggle: todos)
%   \simo{text}{comment}        — inline author comment (toggle: COMMENTS)
%   \rone{revised text}          — blue revision marker (toggle: coloring)
%   \rtwo{revised text}         — violet revision marker (toggle: coloring)
%   \added{text}  \deleted{text}  \replaced{new}{old}  — changes package
%
% FOR FINAL/CAMERA-READY: set all three toggles to false, and switch
%   \usepackage[draft]{changes} to \usepackage[final]{changes}
% ==========================================================================


% --------------------------------------------------------------------------
% UTILITY COMMANDS
% --------------------------------------------------------------------------

% Paint text red — handy for quick emphasis or flagging
% Usage: \red{some important text}
\newcommand{\red}[1]{\textcolor{red}{#1}}


% --------------------------------------------------------------------------
% TODO MARKERS
% --------------------------------------------------------------------------
% Toggle: \TODOStrue / \TODOSfalse

\newif\ifTODOS
\TODOStrue

\newcommand{\todo}[1]{%
  \ifTODOS\textcolor{red}{[TODO: #1]}\fi%
}


% --------------------------------------------------------------------------
% AUTHOR COMMENT COMMANDS
% --------------------------------------------------------------------------
% Toggle: set \COMMENTStrue or \COMMENTSfalse

\newif\ifCOMMENTS
\COMMENTStrue

\ifCOMMENTS
\newcommand{\simo}[2]{\textcolor{red}{#1\bf{ (Simo: #2)}}}
\newcommand{\bob}[2]{\textcolor{blue}{#1\bf{ (Bob: #2)}}}
\newcommand{\alice}[2]{\textcolor{teal}{#1\bf{ (Alice: #2)}}}
\newcommand{\negatron}[2]{\textcolor{purple}{#1\bf{ (Negatron: #2)}}}
\else
\newcommand{\simo}[2]{#1}
\newcommand{\bob}[2]{#1}
\newcommand{\alice}[2]{#1}
\newcommand{\negatron}[2]{#1}
\fi

% Usage: \simo{test sentence}{maybe rephrase this part}
%        \bob{some text}{I disagree with this}
%        \alice{a claim}{needs citation}
%        \negatron{paragraph}{this section is weak}
%
% TO ADD A NEW AUTHOR:
%   1. Pick a color (e.g. orange, brown, cyan — see xcolor docs)
%   2. Add inside BOTH the \ifCOMMENTS and \else branches:
%      \ifCOMMENTS block:
%        \newcommand{\yourname}[2]{\textcolor{yourcolor}{#1\bf{ (Yourname: #2)}}}
%      \else block:
%        \newcommand{\yourname}[2]{#1}
%   3. Optionally add a \definechangesauthor line below for the changes package


% --------------------------------------------------------------------------
% CHANGE TRACKING (changes package)
% --------------------------------------------------------------------------
% https://mirrors.ibiblio.org/CTAN/macros/latex/contrib/changes/changes.english.pdf
\usepackage[draft]{changes}  % shows changes
%\usepackage[final]{changes}  % hides changes

% Define authors for the changes package
\definechangesauthor[name={Simo}, color=red]{SH}
\definechangesauthor[name={Bob}, color=blue]{BO}
\definechangesauthor[name={Alice}, color=teal]{AL}
\definechangesauthor[name={Negatron}, color=purple]{NE}

% Basic usage:
%   \added{This is added text.}
%   \deleted{This text is deleted.}
%   \replaced{new text}{old text}
%
% With author tag:
%   \added[SH]{Author-tagged addition.}
%   \deleted[BO]{Author-tagged deletion.}
%   \replaced[AL]{newer text}{older text}
% --------------------------------------------------------------------------


% --------------------------------------------------------------------------
% REVISION COLORING
% --------------------------------------------------------------------------
% Toggle: \coloringtrue / \coloringfalse
% Use \rone{} for round 1 revisions, \rtwo{} for round 2

\newif\ifcoloring
\coloringtrue

\newcommand{\rone}[1]{%
  \ifcoloring{\color{blue}#1}\else#1\fi%
}

\newcommand{\rtwo}[1]{%
  \ifcoloring{\color{violet}#1}\else#1\fi%
}
 in your preamble (after \usepackage{xcolor})
%   3. Make sure xcolor and changes packages are available
%
% TOGGLES (set to true/false in this file to control visibility):
%   \COMMENTStrue / \COMMENTSfalse    — author comments (\simo, \bob, etc.)
%   \coloringtrue / \coloringfalse    — revision coloring (\rone, \rtwo)
%   \TODOStrue / \TODOSfalse          — todo markers (\todo)
%
% QUICK REFERENCE:
%   \red{text}                  — always-on red text (not toggleable)
%   \todo{remember to fix this} — red TODO text (toggle: todos)
%   \simo{text}{comment}        — inline author comment (toggle: COMMENTS)
%   \rone{revised text}          — blue revision marker (toggle: coloring)
%   \rtwo{revised text}         — violet revision marker (toggle: coloring)
%   \added{text}  \deleted{text}  \replaced{new}{old}  — changes package
%
% FOR FINAL/CAMERA-READY: set all three toggles to false, and switch
%   \usepackage[draft]{changes} to \usepackage[final]{changes}
% ==========================================================================


% --------------------------------------------------------------------------
% UTILITY COMMANDS
% --------------------------------------------------------------------------

% Paint text red — handy for quick emphasis or flagging
% Usage: \red{some important text}
\newcommand{\red}[1]{\textcolor{red}{#1}}


% --------------------------------------------------------------------------
% TODO MARKERS
% --------------------------------------------------------------------------
% Toggle: \TODOStrue / \TODOSfalse

\newif\ifTODOS
\TODOStrue

\newcommand{\todo}[1]{%
  \ifTODOS\textcolor{red}{[TODO: #1]}\fi%
}


% --------------------------------------------------------------------------
% AUTHOR COMMENT COMMANDS
% --------------------------------------------------------------------------
% Toggle: set \COMMENTStrue or \COMMENTSfalse

\newif\ifCOMMENTS
\COMMENTStrue

\ifCOMMENTS
\newcommand{\simo}[2]{\textcolor{red}{#1\bf{ (Simo: #2)}}}
\newcommand{\bob}[2]{\textcolor{blue}{#1\bf{ (Bob: #2)}}}
\newcommand{\alice}[2]{\textcolor{teal}{#1\bf{ (Alice: #2)}}}
\newcommand{\negatron}[2]{\textcolor{purple}{#1\bf{ (Negatron: #2)}}}
\else
\newcommand{\simo}[2]{#1}
\newcommand{\bob}[2]{#1}
\newcommand{\alice}[2]{#1}
\newcommand{\negatron}[2]{#1}
\fi

% Usage: \simo{test sentence}{maybe rephrase this part}
%        \bob{some text}{I disagree with this}
%        \alice{a claim}{needs citation}
%        \negatron{paragraph}{this section is weak}
%
% TO ADD A NEW AUTHOR:
%   1. Pick a color (e.g. orange, brown, cyan — see xcolor docs)
%   2. Add inside BOTH the \ifCOMMENTS and \else branches:
%      \ifCOMMENTS block:
%        \newcommand{\yourname}[2]{\textcolor{yourcolor}{#1\bf{ (Yourname: #2)}}}
%      \else block:
%        \newcommand{\yourname}[2]{#1}
%   3. Optionally add a \definechangesauthor line below for the changes package


% --------------------------------------------------------------------------
% CHANGE TRACKING (changes package)
% --------------------------------------------------------------------------
% https://mirrors.ibiblio.org/CTAN/macros/latex/contrib/changes/changes.english.pdf
\usepackage[draft]{changes}  % shows changes
%\usepackage[final]{changes}  % hides changes

% Define authors for the changes package
\definechangesauthor[name={Simo}, color=red]{SH}
\definechangesauthor[name={Bob}, color=blue]{BO}
\definechangesauthor[name={Alice}, color=teal]{AL}
\definechangesauthor[name={Negatron}, color=purple]{NE}

% Basic usage:
%   \added{This is added text.}
%   \deleted{This text is deleted.}
%   \replaced{new text}{old text}
%
% With author tag:
%   \added[SH]{Author-tagged addition.}
%   \deleted[BO]{Author-tagged deletion.}
%   \replaced[AL]{newer text}{older text}
% --------------------------------------------------------------------------


% --------------------------------------------------------------------------
% REVISION COLORING
% --------------------------------------------------------------------------
% Toggle: \coloringtrue / \coloringfalse
% Use \rone{} for round 1 revisions, \rtwo{} for round 2

\newif\ifcoloring
\coloringtrue

\newcommand{\rone}[1]{%
  \ifcoloring{\color{blue}#1}\else#1\fi%
}

\newcommand{\rtwo}[1]{%
  \ifcoloring{\color{violet}#1}\else#1\fi%
}
 in your preamble (after \usepackage{xcolor})
%   3. Make sure xcolor and changes packages are available
%
% TOGGLES (set to true/false in this file to control visibility):
%   \COMMENTStrue / \COMMENTSfalse    — author comments (\simo, \bob, etc.)
%   \coloringtrue / \coloringfalse    — revision coloring (\rone, \rtwo)
%   \TODOStrue / \TODOSfalse          — todo markers (\todo)
%
% QUICK REFERENCE:
%   \red{text}                  — always-on red text (not toggleable)
%   \todo{remember to fix this} — red TODO text (toggle: todos)
%   \simo{text}{comment}        — inline author comment (toggle: COMMENTS)
%   \rone{revised text}          — blue revision marker (toggle: coloring)
%   \rtwo{revised text}         — violet revision marker (toggle: coloring)
%   \added{text}  \deleted{text}  \replaced{new}{old}  — changes package
%
% FOR FINAL/CAMERA-READY: set all three toggles to false, and switch
%   \usepackage[draft]{changes} to \usepackage[final]{changes}
% ==========================================================================


% --------------------------------------------------------------------------
% UTILITY COMMANDS
% --------------------------------------------------------------------------

% Paint text red — handy for quick emphasis or flagging
% Usage: \red{some important text}
\newcommand{\red}[1]{\textcolor{red}{#1}}


% --------------------------------------------------------------------------
% TODO MARKERS
% --------------------------------------------------------------------------
% Toggle: \TODOStrue / \TODOSfalse

\newif\ifTODOS
\TODOStrue

\newcommand{\todo}[1]{%
  \ifTODOS\textcolor{red}{[TODO: #1]}\fi%
}


% --------------------------------------------------------------------------
% AUTHOR COMMENT COMMANDS
% --------------------------------------------------------------------------
% Toggle: set \COMMENTStrue or \COMMENTSfalse

\newif\ifCOMMENTS
\COMMENTStrue

\ifCOMMENTS
\newcommand{\simo}[2]{\textcolor{red}{#1\bf{ (Simo: #2)}}}
\newcommand{\bob}[2]{\textcolor{blue}{#1\bf{ (Bob: #2)}}}
\newcommand{\alice}[2]{\textcolor{teal}{#1\bf{ (Alice: #2)}}}
\newcommand{\negatron}[2]{\textcolor{purple}{#1\bf{ (Negatron: #2)}}}
\else
\newcommand{\simo}[2]{#1}
\newcommand{\bob}[2]{#1}
\newcommand{\alice}[2]{#1}
\newcommand{\negatron}[2]{#1}
\fi

% Usage: \simo{test sentence}{maybe rephrase this part}
%        \bob{some text}{I disagree with this}
%        \alice{a claim}{needs citation}
%        \negatron{paragraph}{this section is weak}
%
% TO ADD A NEW AUTHOR:
%   1. Pick a color (e.g. orange, brown, cyan — see xcolor docs)
%   2. Add inside BOTH the \ifCOMMENTS and \else branches:
%      \ifCOMMENTS block:
%        \newcommand{\yourname}[2]{\textcolor{yourcolor}{#1\bf{ (Yourname: #2)}}}
%      \else block:
%        \newcommand{\yourname}[2]{#1}
%   3. Optionally add a \definechangesauthor line below for the changes package


% --------------------------------------------------------------------------
% CHANGE TRACKING (changes package)
% --------------------------------------------------------------------------
% https://mirrors.ibiblio.org/CTAN/macros/latex/contrib/changes/changes.english.pdf
\usepackage[draft]{changes}  % shows changes
%\usepackage[final]{changes}  % hides changes

% Define authors for the changes package
\definechangesauthor[name={Simo}, color=red]{SH}
\definechangesauthor[name={Bob}, color=blue]{BO}
\definechangesauthor[name={Alice}, color=teal]{AL}
\definechangesauthor[name={Negatron}, color=purple]{NE}

% Basic usage:
%   \added{This is added text.}
%   \deleted{This text is deleted.}
%   \replaced{new text}{old text}
%
% With author tag:
%   \added[SH]{Author-tagged addition.}
%   \deleted[BO]{Author-tagged deletion.}
%   \replaced[AL]{newer text}{older text}
% --------------------------------------------------------------------------


% --------------------------------------------------------------------------
% REVISION COLORING
% --------------------------------------------------------------------------
% Toggle: \coloringtrue / \coloringfalse
% Use \rone{} for round 1 revisions, \rtwo{} for round 2

\newif\ifcoloring
\coloringtrue

\newcommand{\rone}[1]{%
  \ifcoloring{\color{blue}#1}\else#1\fi%
}

\newcommand{\rtwo}[1]{%
  \ifcoloring{\color{violet}#1}\else#1\fi%
}
